(i) For a quasi-free state the four-point function of the $\alpha$'s is the Pfaffian (Wick) sum of products of two-point functions (Fact \claimref{hawking.imported-quasi-free-states-car-algebra}). For $O = \tfrac12\sum_{IJ}F_{IJ}\alpha_{I}^{\dagger}\alpha_{J}$ the connected part of $\omega(OO)$ keeps the two contractions that join the two factors; with $\Gamma_{IJ} = \omega(\alpha_{J}^{\dagger}\alpha_{I})$ and $\omega(\alpha_{I}\alpha_{J}^{\dagger}) = (1 - \Gamma)_{IJ}$, and the particle-hole form of $F$ (which makes the two contractions equal), they sum to $\tfrac12\Tr(F\Gamma F(1 - \Gamma))$. When $\omega$ is pure, $\Gamma$ is a projector, $\Tr(F\Gamma F(1 - \Gamma)) = \|(1 - \Gamma)F\Gamma\|_{2}^{2}$, and $\|[F, \Gamma]\|_{2}^{2} = \|(1 - \Gamma)F\Gamma\|_{2}^{2} + \|\Gamma F(1 - \Gamma)\|_{2}^{2} = 2\|(1 - \Gamma)F\Gamma\|_{2}^{2}$, so $\mathrm{Var} = \tfrac14\|[F, \Gamma]\|_{2}^{2}$. In the gauge-invariant case $\Gamma = E \oplus (1 - E^{T})$ and $F = f \oplus (-f^{T})$, the two blocks contribute equally, and $\mathrm{Var}(N_{f}) = \Tr(fEf(1 - E))$. The check computes both sides on the Fock space: for a mixed gauge-invariant state built as a product over rotated modes with rational orthogonal rotation, for a projector, for a mixed state with pairing obtained by restricting a pure Thouless state of five modes to three, and for the stored pure pair state.

(ii) Let $\widetilde\omega$ be a pure quasi-free state of a doubled algebra (system and partner modes) whose restriction to the system is $\omega$, with projector $\widetilde\Gamma$ whose system block is $\Gamma$. Such a purification exists for every quasi-free state (standard, Araki; from memory); in the gauge-invariant case it is explicit, $\widetilde E = \bigl(\begin{smallmatrix}E & S \\ S & 1 - E\end{smallmatrix}\bigr)$ with $S = (E(1 - E))^{1/2}$, which is a projector because $S$ commutes with $E$ and $S^{2} = E(1 - E)$. Extend $O$ by zero on the partner, $\widetilde F = F \oplus 0$, so $\mathrm{Var}_{\widetilde\omega}(\widetilde O) = \mathrm{Var}_{\omega}(O)$. Let $\widetilde\Gamma_{0} = \Gamma_{0} \oplus \Gamma_{0}'$ with $\Gamma_{0}'$ any projector of the partner's Nambu form (in the gauge-invariant case $1 - E_{0}$), so that $\mathrm{Var}_{\widetilde\omega_{0}}(\widetilde O) = \mathrm{Var}_{\omega_{0}}(O)$, and $\Delta = \widetilde\Gamma - \widetilde\Gamma_{0}$. By (i) for pure states and the triangle inequality of the Hilbert-Schmidt norm,
\[
2\,\mathrm{sd}_{\omega}(O) = \|[\widetilde F, \widetilde\Gamma]\|_{2} \le \|[\widetilde F, \widetilde\Gamma_{0}]\|_{2} + \|\widetilde F\Delta - \Delta\widetilde F\|_{2} \le 2\,\mathrm{sd}_{\omega_{0}}(O) + 2\|\widetilde F\Delta\|_{2},
\]
and $\|\widetilde F\Delta\|_{2}^{2} = \Tr(\widetilde F^{2}\Delta^{2}) = \Tr(F^{2}X)$ with $X$ the system block of $\Delta^{2}$. That block is $(\Gamma - \Gamma_{0})^{2} + \widetilde\Gamma_{12}\widetilde\Gamma_{21}$, and the system block of $\widetilde\Gamma^{2} = \widetilde\Gamma$ reads $\Gamma^{2} + \widetilde\Gamma_{12}\widetilde\Gamma_{21} = \Gamma$; hence $X = (\Gamma - \Gamma_{0})^{2} + \Gamma(1 - \Gamma)$, independent of the purification and of $\Gamma_{0}'$, and $X \ge 0$ as a sum of two positive matrices. Its trace is $\Tr(\Gamma - 2\Gamma\Gamma_{0} + \Gamma_{0}) = \Tr(\Gamma(1 - \Gamma_{0})) + \Tr((1 - \Gamma)\Gamma_{0})$; in the quasi-particle frame of $\omega_{0}$, where $\Gamma_{0} = 0 \oplus 1$, both terms equal $\omega(N_{0})$, and the expression is invariant under the Bogoliubov change of frame, so $\Tr X = 2\omega(N_{0})$; the operator $\tfrac12\alpha^{\dagger}(1 - 2\Gamma_{0})\alpha + \tfrac n2$ is $N_{0}$ in that frame. The gauge-invariant form follows with $F = f \oplus (-f^{T})$: $\Tr(F^{2}X) = 2\Tr(f^{2}X_{E})$. The check verifies on the Fock space, with exact rationals: the purification is a projector restricting to $E$; $\|\widetilde E - \widetilde E_{0}\|_{2}^{2} = 2\Tr X_{E}$; the system block of $\Delta^{2}$ is $X_{E}$; $\|\widetilde F\Delta\|_{2}^{2} = \Tr(f^{2}X_{E})$; $2\mathrm{Var} = \|[\widetilde f, \widetilde E]\|_{2}^{2}$; the inequality for several $f$; in the pairing case the identity $\widetilde\Gamma_{12}\widetilde\Gamma_{21} = \Gamma(1 - \Gamma)$ on the five-mode Thouless purification, $\omega_{0}(N_{0}) = 0$, $\Tr X = 2\omega(N_{0})$ and the inequality, including the pure pair state $\tfrac35\ket{00} + \tfrac45\ket{11}$ of Proposition \claimref{synthesis.quadratic-flows-cannot-thermalise-descent}(iii) with the observable $a^{\dagger}b + b^{\dagger}a$ coupling mode and partner, against the empty reference. Quasi-freeness is used in (i) and in the purification: for the mixture $\tfrac12(\ket{0}\bra{0} + \ket{1\dots1}\bra{1\dots1})$ of five modes, total number $N$, $\mathrm{Var} = 25/4$ while $E = \tfrac12$, $X_{E} = \tfrac12$, and the bound would give $2\Tr X_{E} = 5$; the check records the violation.

(iii) The descent is the one-way conditional expectation of Fact \claimref{singularity.imported-descent-one-way-conditional-expectation}; $E_{\mathrm{d}}$ is the average over $\{\mathrm{id}, \mathrm{Ad}\,\sigma\}$ (Theorem \claimref{singularity.descent-parity-average}), so $\omega\circ E_{\mathrm{d}} = \tfrac12(\omega + \omega\circ\mathrm{Ad}\,\sigma)$, and for any two states the variance of a mixture is the mean of the variances plus the variance of the means: $\mathrm{Var}_{\omega\circ E_{\mathrm{d}}}(O) = \tfrac12\mathrm{Var}_{\omega}(O) + \tfrac12\mathrm{Var}_{\omega}(\sigma O\sigma) + \tfrac14(\omega(O) - \omega(\sigma O\sigma))^{2}$. Since $\sigma a_{0}\sigma = -a_{0}$ and $\sigma a_{j}\sigma = a_{j}$ for $j \ne 0$, $\sigma\alpha\sigma = S\alpha$ and $\sigma O\sigma = \tfrac12\alpha^{\dagger}SFS\alpha$, a quadratic element, and $\omega\circ\mathrm{Ad}\,\sigma$ is quasi-free with $\Gamma \mapsto S\Gamma S$, so (ii) applies to both variances. $F - SFS$ is twice the part of $F$ with exactly one index in $K$, so $D = \tfrac12\Tr((F - SFS)\Gamma) = \sum_{I \in K, J \notin K}(F_{IJ}\Gamma_{JI} + F_{JI}\Gamma_{IJ})$, and by Cauchy-Schwarz with $\|\Gamma e_{I}\| \le 1$ ($0 \le \Gamma \le 1$), $|D| \le 2\sum_{I \in K}r_{I}$, $r_{I}^{2} = \sum_{J \notin K}|F_{IJ}|^{2}$, so $D^{2} \le 8(r_{1}^{2} + r_{2}^{2})$. Both identities and the bound are checked on the Fock space for the restricted Thouless state.

(iv) Take a comoving box of volume $V$ and the cosine moment $C = \int_{V}\cos(\vec k\cdot\vec x)(\rho - \bar\rho)$ of the energy density of the fermion modes, a one-body observable $N_{f_{C}}$ with $f_{C} = \tfrac12(cH + Hc)$ up to a constant, $c$ the multiplication by $\cos(\vec k\cdot\vec x)$ and $H$ the one-particle Hamiltonian; likewise the sine moment $S$. With $A = cH$, $f_{C}^{2} = \tfrac14(A + A^{\dagger})^{2} \le \tfrac12(AA^{\dagger} + A^{\dagger}A) = \tfrac12(cH^{2}c + Hc^{2}H) \le \tfrac12(cH^{2}c + H^{2})$. For a homogeneous gauge-invariant deposit $X_{E}$ commutes with translations, $X_{E} = \bigoplus_{p}X_{p}$ on the band (spin and charge) space at momentum $p$, with $x_{p} = \Tr X_{p}$; $H^{2} = \omega_{p}^{2}$ there; the momentum-diagonal part of $cH^{2}c$ is $\tfrac14(\omega_{p + k}^{2} + \omega_{p - k}^{2})$, and $\omega_{p \pm k} \le \omega_{p} + k$ (the dispersion is Lipschitz with constant one). Hence $\Tr(f_{C}^{2}X_{E}) \le \tfrac34\sum_{p}(\omega_{p} + k)^{2}x_{p} = \tfrac34 V\bar n\langle(\omega + k)^{2}\rangle$, with $\bar n = V^{-1}\sum_{p}x_{p}$ the density of particles plus antiparticles, and $\bar\rho = \bar n\langle\omega\rangle$ above the vacuum. By (ii), $\mathrm{Var}(C)^{1/2} \le \mathrm{Var}_{0}(C)^{1/2} + (\tfrac32 V\bar n\langle(\omega + k)^{2}\rangle)^{1/2}$, and the same for $S$. The spectrum of the contrast is $P(k) = (\mathrm{Var}\,C + \mathrm{Var}\,S)/(\bar\rho^{2}V)$ (the energy densities at different points commute), so by the triangle inequality in $\mathbb{R}^{2}$, $P^{1/2} \le P_{\mathrm{vac}}^{1/2} + (3c_{k}/\bar n)^{1/2}$, and $\Delta^{2} = k^{3}P/2\pi^{2}$ gives the stated bound.

Numbers. For massless Fermi-Dirac quanta (the I-state, Theorem \claimref{hawking.i-state}), $\langle\omega^{m}\rangle = \Gamma(m + 3)\eta(m + 3)T^{m}/(\Gamma(3)\eta(3))$ with $\eta(s) = (1 - 2^{1 - s})\zeta(s)$ (standard), so $c = \langle\omega^{2}\rangle/\langle\omega\rangle^{2} = \tfrac43\eta(5)\eta(3)/\eta(4)^{2}$; with $\zeta(3) < 1.203$, $\zeta(5) < 1.04$, $\zeta(4) = \pi^{4}/90$ and $\pi > 333/106$, $c < 1.309$, and $c_{k} \le c(1 + k/\langle\omega\rangle)^{2}$ with $\langle\omega\rangle > 3T$. The record fixes $T_{0} = 2.7255$ K and $k_{*} = 0.05\ \mathrm{Mpc}^{-1}$ (Fact \claimref{synthesis.imported-planck-2018-lcdm-parameters}), so $k_{*}/(a_{0}T_{0}) \in (1.35, 1.37) \times 10^{-27}$ (with $k_{B}$, $\hbar c$ and the megaparsec at their defined or CODATA values). Entropy conservation from the descent to today (standard; the stored finite-range bound assumes it too) gives $(a_{d}T_{d})^{3} = (a_{0}T_{0})^{3}g_{s0}/g_{sd}$ with $g_{s0} = 43/11$; for $g_{sd} \le 10^{4}$, $k_{*}/(a_{d}T_{d}) < 14 \times 1.37 \times 10^{-27}$, so $c_{k_{*}} < 1.31$. One Weyl species (particle and antiparticle) in the I-state has $\bar n = \tfrac{3\zeta(3)}{2\pi^{2}}(a_{d}T_{d})^{3}$ comoving, and more species only raise $\bar n$. With $\zeta(3) > 6/5$ and $333/106 < \pi < 355/113$ the shot term $3c_{k}k_{*}^{3}/(2\pi^{2}\bar n)$ is below $10^{-77}$ (about $7 \times 10^{-78}$ at $g_{sd} = 10^{4}$, about $7 \times 10^{-80}$ at $106.75$), and $10^{68}$ times it is below $A_{s} \ge 2.07 \times 10^{-9}$ (the comparison with $A_{s}$ is made under hypothesis (FR) of the formal note on the conversion). $T_{d}$ drops out: only the comoving number of quanta enters, and it is the record's. The check reproduces every bracket in exact arithmetic.

Consequences. The pure pair deposits of Proposition \claimref{synthesis.quadratic-flows-cannot-thermalise-descent}(iii), with partners placed anywhere on the descent slice, are pure quasi-free states of the slice's CAR algebra; the I-state is quasi-free and gauge-invariant (Theorem \claimref{hawking.i-state}); restriction to a subalgebra and quadratic flows keep quasi-freeness. None of these hypotheses involves the correlation range or the partner placement, so (ii) and (iv) bound every such deposit by its vacuum plus the shot term $(3c_{k}k^{3}/2\pi^{2}\bar n)^{1/2}$, and (iii) adds a region-independent constant for the descent's conditional expectation. The shot term depends on the occupation profile through $\bar n$ and $c_{k}$: the numbers above hold for the I-state and for any profile with at least one thermal Weyl species' comoving number of quanta and $c_{k} \le 1.31$. They do not hold for every profile: a gauge-invariant product state with of order one quantum per pivot volume has $\bar n$ of order $k_{*}^{3}$ and a shot term of order one, and such a deposit reaches today's radiation only through entropy production, so the record does not fix its number. Since the energy density of the fermion modes is part of the total, $\delta\rho_{F}/\bar\rho_{\mathrm{tot}} \le \delta\rho_{F}/\bar\rho_{F}$, and the bound holds for their contribution to the total contrast; other components are not bounded here. The vacuum's own term is the descended vacuum's spectrum: for massless fermions the conformal vacuum, whose energy-density correlator is homogeneous of degree $-8$, so $\Delta_{\delta, \mathrm{vac}}^{2}$ scales as $(k/a_{d}T_{d})^{8}$ against the thermal density (argued, standard), negligible at the pivot; the frozen-point reading of Proposition \claimref{synthesis.frozen-point-geometry-reads-descent-as-vacuum} gives index $2$, not the measured index, for the vacuum-like class, and the finite-range bound (Proposition \claimref{synthesis.finite-range-descent-deposit-cannot-seed-cmb}) gives index $3$ for white noise, which is what the shot term is (under (FR)). Not covered: states that are not quasi-free (mixtures over a classical modulating field, where the variance of the local means adds without bound, and entangled non-Gaussian states, as the cat mixture shows), and bosonic quasi-free states, for which $\mathrm{Var}(N_{f}) = \Tr(fEf(1 + E))$ grows with the occupations.

The conversion note (version 5). The comparisons with $A_{s}$ and the index $3$ in the paragraphs above are made under (FR). The constraint and the amplitude $\mathcal{R}_{\infty} = \psi + \tfrac14\delta - 2Hq/(\rho + p)$ of the note are those of Proposition \claimref{synthesis.descent-slice-constraints-and-the-three-readings-of-a-white-deposit}, the result on the descent slice's constraints; they are cited, not re-derived. As evidence for this amendment the check re-verifies the two coefficients the note uses, by exact series in $x = k\eta = 1/X$: for the radiation-era equation $\Phi'' + \tfrac4x\Phi' + \tfrac13\Phi = 0$, with $\mathcal{R} = \Phi + \tfrac12(x\Phi' + \Phi)$ (which is $\tfrac32\Phi$ on the growing solution at leading order, Fact \claimref{synthesis.imported-large-scale-potential-and-sachs-wolfe-relations} at $w = \tfrac13$) and the comoving contrast $\Delta_{c} = -\tfrac23x^{2}\Phi$, data $(\mathcal{R}, \Delta_{c})$ at $x$ have the growing amplitude $\mathcal{R}_{\infty} = \mathcal{R}\,(1 + \tfrac{x^{2}}{6} + \dots) + \Delta_{c}\,(\tfrac14 - \tfrac{x^{2}}{24} + \dots)$, and on a slice $\mathcal{R} = \psi - Hq/(\rho + p)$, $\Delta_{c} = \delta - 3Hq/\rho$. With $\psi = -\tfrac32X^{2}\delta$ (the constraint at $\kappa_{\theta} = 0$, taken from that Proposition and not tested here) the terms of order $\delta$ cancel, $\mathcal{R}_{\infty} = -\tfrac32X^{2}\delta - 2Hq/(\rho + p) + O(\delta/X^{2})$. In the convention of Fact \claimref{synthesis.imported-superhorizon-zeta-conservation}, $-\zeta = \mathcal{R} - \tfrac14\Delta_{c}$ on a slice, which differs from $\mathcal{R}$ by a term of order $x^{2}\Phi$, the order neglected there; under (FR) the slice has $\mathcal{R} = 0$ and $\Delta_{c} = \delta$, so $\zeta = +\tfrac14\delta$ on the slice and $\zeta_{\infty} = -\mathcal{R}_{\infty} = -\tfrac14\delta$ afterwards. The thermal ratio is this entity's own: the variance formula of (i) gives $\int\omega^{2}n(1 - n) = T^{2}\partial\rho/\partial T = 4\rho T$ for the energy form and a third of it, $(\rho + p)T$, for the longitudinal momentum form, at $k \to 0$, with no cross term by isotropy; hence $P_{\delta} = 4T/\rho$, $P_{v} = T/(\rho + p) = \tfrac{3}{16}P_{\delta}$ and $4X^{2}P_{v} = 12X^{2}\cdot\tfrac{1}{16}P_{\delta}$. With $\Delta_{\delta}^{2} = \tfrac{60}{\pi^{4}g}(k_{d}/T_{d})^{3}$ at the physical wavenumber $k_{d} = H_{d}/X$, the reading (UE) gives $\tfrac{135}{\pi^{4}g}(H_{d}/T_{d})^{3}X$, times the share $f$ for the fermion modes alone. Numbers, each an exact bracket in the check: with $H_{d} = \pi\sqrt{g/90}\,T_{d}^{2}/\bar M_{P}$ the square of the (UE) amplitude is $(135f)^{2}g\,(T_{d}/\bar M_{P})^{6}X^{2}/(90^{3}\pi^{2})$, between $5.7^{2}$ and $5.8^{2}$ at the pivot for $T_{d} = 10^{12}$ GeV, $g = 106.75$, $f = \tfrac12$ and the stored lower bound $X = 3.2 \times 10^{20}$, and between $(26A_{s})^{2}$ and $(28A_{s})^{2}$ at $10^{10}$ GeV where $X$ is a hundred times smaller; under (FM) the shot term of (iv) at $g = 106.75$, below $8 \times 10^{-80}$, times $(1 + 12X^{2})/16$ with $X < 3.44 \times 10^{20}$, is below $7.1 \times 10^{-39}$.
